% ===========================================================================
% Section II-B -- Pre-Processing and Kurtosis-Guided Bandpass Filter Design
% Owner: Member 2
% ===========================================================================
\subsection{Pre-Processing and Kurtosis-Guided Bandpass Filter Design}

\subsubsection{DC Offset Removal and Detrending}
Sensor electronics add a constant bias, which appears as a large peak at
$0$~Hz whose sidelobes leak into the low-frequency bins that matter. Before any
filtering, the signal $x[n]$ of length $N$ is centered by subtracting its sample
mean:
\begin{equation}
  x_0[n] = x[n] - \mu_x, \quad \text{where} \quad
  \mu_x = \frac{1}{N}\sum_{m=0}^{N-1} x[m]
\end{equation}
In the DFT sense this sets $X_0(e^{j0})=\sum_{n=0}^{N-1}x_0[n]=0$ and leaves
every other coefficient untouched, so the filter bank sees only the vibration
riding on a zero baseline.

\subsubsection{Sub-Band Kurtosis Resonance Selection}
Below about $1000$~Hz the spectrum is crowded with shaft harmonics, structural
vibration and load-dependent friction noise, which hide the bearing impulses.
Each time a rolling element hits a defect it excites a high-frequency resonance
of the housing, and that band carries sharp bursts while the low-frequency
interference stays stationary. The resonance belongs to the structure, not to
the fault, which is why one band collects many bursts at once.

The resonance depends on the machine and its mounting, so a fixed bandpass
filter cannot find it reliably and the band has to be picked per record. A
filter bank built on the Fast Kurtogram~\cite{antoni2007} does this without
training data. The Nyquist interval $[0, f_s/2]$ is split at
dyadic and half-integer levels. On a $12$~kHz record the limit is $6000$~Hz, so
the dyadic levels give bands of $3000$, $1500$, $750$~Hz and so on, while the
half-integer levels add widths of $2000$, $1000$ and $545.5$~Hz. The $545.5$~Hz
band used in Section~IV comes from that set. At level $k$ there are
$M_k = \text{round}(2^k)$ bands, and
band $i$ has bandwidth and center frequency
\begin{equation}
  \Delta f_k = \frac{f_s}{2 M_k}, \quad
  f_{c,i} = \left(i + \frac{1}{2}\right)\Delta f_k
\end{equation}

The search is restricted to $[1000~\text{Hz},\ f_s/2-150~\text{Hz}]$: the lower
guard skips the crowded region below $1000$~Hz and the upper guard keeps the
acquisition anti-alias transition band out of the search. Any sub-band crossing
either guard is scored $-\infty$, so an excluded band can never win.

For each band output $y_{k,i}[n]$, impulsiveness is measured by the fourth
standardized central moment, the Kurtosis~\cite{antoni2006}:
\begin{equation}
  K(f_{c,i}, \Delta f_k) = \frac{\frac{1}{N}\sum_{n=0}^{N-1}
  \left(y_{k,i}[n] - \mu_{y}\right)^4}{\left(\frac{1}{N}\sum_{n=0}^{N-1}
  \left(y_{k,i}[n] - \mu_{y}\right)^2\right)^2}
\end{equation}
with $\mu_y$ the mean of the band signal. The fourth power lets one large
transient outweigh many small samples, while the squared variance $\sigma_y^4$
normalizes out the band energy so a loud band does not win by being loud.
Gaussian noise gives $K \approx 3$; a band of repeated impacts gives $K$ well
above $3$. The band with the largest $K$ is selected:
\begin{equation}
  \{f_c^*, BW^*\} = \arg\max_{f_{c,i}, \, \Delta f_k} K(f_{c,i}, \Delta f_k)
\end{equation}
Its cutoffs are $f_{\text{low}} = f_c^* - BW^*/2$ and
$f_{\text{high}} = f_c^* + BW^*/2$, that is
$\omega_1 = 2\pi f_{\text{low}}/f_s$ and $\omega_2 = 2\pi f_{\text{high}}/f_s$.

\subsubsection{Bandpass Filter Design: Kaiser FIR versus Butterworth IIR}
Two filters are designed for that passband, a linear-phase FIR and a
Butterworth IIR, both rejecting out-of-band motor interference and broadband
noise. The FIR is Type-I, with even order $M$ and $L = M + 1$ taps: the ideal
bandpass impulse response is truncated around its symmetry center $\alpha = M/2$
and multiplied by a Kaiser window~\cite{oppenheim}:
\begin{equation}
\begin{split}
  h_{\text{FIR}}[n] &= \left( \frac{\sin(\omega_2(n - \alpha)) -
  \sin(\omega_1(n - \alpha))}{\pi(n - \alpha)} \right) \\
  &\quad \times \frac{I_0\!\left(\beta \sqrt{1 -
  \left(\frac{n - \alpha}{\alpha}\right)^2}\right)}{I_0(\beta)}
\end{split}
\end{equation}
for $0 \le n \le M$. $I_0(\cdot)$ is the zeroth-order modified Bessel function
of the first kind, and $\beta$ sets the trade between transition width and
sidelobe attenuation. The taps are symmetric,
$h_{\text{FIR}}[n] = h_{\text{FIR}}[M - n]$, so the phase is exactly linear and
the delay is flat:
\begin{equation}
  \tau_{g,\text{FIR}}(\omega) =
  -\frac{d\angle H_{\text{FIR}}(e^{j\omega})}{d\omega} = \frac{M}{2}
  \text{ samples}
\end{equation}
An FIR of this form has all its poles at $z = 0$, so it is stable whatever the
coefficients do; the price is arithmetic, since every extra tap buys a steeper
transition at one more operation per sample. Section~IV uses $\beta = 8$, which
targets roughly $-80$~dB of stopband attenuation.

The second design is a $4^{\text{th}}$-order Butterworth bandpass, obtained with
the bilinear transform~\cite{proakis}. That transform maps the left half of the $s$-plane into
the unit circle:
\begin{equation}
  s = \frac{2}{T_s}\left(\frac{1 - z^{-1}}{1 + z^{-1}}\right)
\end{equation}
with $T_s = 1/f_s$. The mapping compresses frequency non-linearly, so the
digital cutoffs are pre-warped before the prototype is transformed:
\begin{equation}
  \Omega_k = \frac{2}{T_s} \tan\!\left(\frac{\omega_k}{2}\right),
  \quad k \in \{1, 2\}
\end{equation}
After pre-warping the $-3$~dB points land exactly on $f_{\text{low}}$ and
$f_{\text{high}}$ with a maximally flat passband. Those are single-pass points:
the forward--backward filter squares the magnitude response, so they appear at
$-6$~dB in the effective zero-phase response. It runs on buffered $2.0$~s
blocks, so the whole block is available before the reverse pass begins.

\subsubsection{Cascade SOS Realization and Zero-Phase Filtering}
A $4^{\text{th}}$-order bandpass has eight poles, and in direct form a small
coefficient error moves them far enough to cross the unit circle. The fix is to
factor the transfer function into $K = 4$ second-order sections:
\begin{equation}
  H_{\text{IIR}}(z) = G \prod_{k=1}^{K} H_k(z) = G \prod_{k=1}^{K}
  \frac{b_{0k} + b_{1k}z^{-1} + b_{2k}z^{-2}}
       {1 + a_{1k}z^{-1} + a_{2k}z^{-2}}
\end{equation}
where $G$ is the overall gain. Each section holds one conjugate pair
$p_k, p_k^* = r_k e^{\pm j\theta_k}$. Stability needs the ROC $|z| > r_{\max}$
to contain the unit circle, so every pole radius must satisfy
\begin{equation}
  r_k = \sqrt{a_{2k}} < 1, \quad \forall k \in \{1, 2, \dots, K\}
\end{equation}

The cascade is cheap, but its phase is not linear near the band edges, so the
frequency-dependent group delay smears a short impact and lowers its peak.
Forward--backward filtering removes that: the record is filtered, reversed,
filtered again, and reversed back,
\begin{equation}
  Y_{\text{BP}}(z) = H(z) H(z^{-1}) X_0(z)
\end{equation}
On the unit circle, $H(e^{-j\omega}) = H^*(e^{j\omega})$ for real coefficients,
so
\begin{equation}
  H_{\text{zero-phase}}(e^{j\omega}) = H(e^{j\omega}) H^*(e^{j\omega}) =
  |H(e^{j\omega})|^2
\end{equation}
which is real and non-negative, so the phase and the group delay are zero. The
impacts keep their time positions and shape, and the stopband attenuation
doubles in decibels: $-40$~dB becomes $-80$~dB. The passband edge shifts a
little inward, because the effective filter is the square of the designed
one.
